\begin{helper}
Fix a terminal coordinate set \(U\), support labels \(B,J\), branch labels
\(\beta\), and transcript labels \(\tau=(P_\tau,A_\tau,Z_\tau)\).  Let
\(\operatorname{blueTrace}(\beta)\) be the complete terminal blue trace of
\(\beta\), let \(\operatorname{redSupport}(\beta,J)\) be the fixed red
support for \(J\), let \(\mathsf{Realized}(B,J,\beta,\tau)\) be the realized
cell predicate, and let \(\mathsf{Comp}(A;B,J,\beta,\tau)\) be compatibility
of a complete terminal assignment \(A\).  Also fix a deterministic transcript
scan
\[
\operatorname{scan}:2^U\to \mathcal P\cup\{\bot\}.
\]

Assume branch labels are functional for complete terminal blue truth tables.
Assume every realized \((\beta,\tau)\) has the exported geometry
\[
\operatorname{redSupport}(\beta,J)\subseteq U,\qquad
\operatorname{redSupport}(\beta,J)\hbox{ is red-only},
\]
\[
P_\tau,A_\tau\subseteq U,\qquad
P_\tau\cap A_\tau=\emptyset,\qquad
B\cup\operatorname{redSupport}(\beta,J)\subseteq P_\tau,
\]
\[
e\in\operatorname{blueTrace}(\beta)\Longleftrightarrow e\in P_\tau,
\qquad
e\notin\operatorname{blueTrace}(\beta)\Longleftrightarrow e\in A_\tau
\]
for every blue coordinate \(e\in U\), and \(Z_\tau\subseteq A_\tau\).
Assume compatible complete assignments contain \(P_\tau\), avoid
\(A_\tau\setminus Z_\tau\), and recover the complete blue trace of
\(\beta\).

Assume the actual fixed-\((B,J)\) scan interface.  First, residual scan
stability holds: if \(A_{\rm res}\subseteq U\) contains
\[
P_\tau\setminus(B\cup\operatorname{redSupport}(\beta,J))
\]
and avoids \(A_\tau\setminus Z_\tau\), then
\[
\operatorname{scan}\bigl((A_{\rm res}\setminus Z_\tau)\cup
  (B\cup\operatorname{redSupport}(\beta,J))\bigr)=\tau .
\]
Second, scan recognition holds from the full forced scan data: if a
realized \((\beta,\tau)\) has
\[
P_\tau\subseteq A,\qquad A_\tau\cap A=\emptyset,
\]
and the complete blue trace of \(\beta\) agrees with \(A\) on all blue
coordinates of \(U\), then \(\operatorname{scan}(A)=\tau\).  Third, the fixed residual
first-mismatch or prefix-pruning uniqueness input holds: if one residual
assignment lies in the residual cylinders of two realized pairs and the two
delete-and-reinsert complete assignments are compatible with those pairs,
then the two branches and transcripts are equal.

Then the two residual outputs used by the fixed-\((B,J)\) support tree hold.
First, under the residual span and avoidance conditions for
\((\beta,\tau)\), put
\[
A^\ast=(A_{\rm res}\setminus Z_\tau)\cup
  (B\cup\operatorname{redSupport}(\beta,J))
\]
Then any realized transcript \((\beta',\tau')\) whose full forced scan data
are satisfied by \(A^\ast\),
\[
P_{\tau'}\subseteq A^\ast,\qquad A_{\tau'}\cap A^\ast=\emptyset,
\]
and whose complete blue trace agrees with \(A^\ast\), is equal to
\(\tau\).  Second, if the same residual assignment reconstructs compatibly
to two realized pairs \((\beta,\tau)\) and \((\beta',\tau')\), then
\(\beta=\beta'\) and \(\tau=\tau'\).
\end{helper}
\begin{proof}
The point is that transcript equality is now derived from deterministic scan
output, not from realized geometry alone.  The geometry and forward
compatibility assumptions are retained only to translate compatibility
certificates back into residual-cylinder span and avoidance conditions.

Fix \(A_{\rm res}\subseteq U\) in the residual cylinder of a realized
\((\beta,\tau)\), and put
\[
A^\ast=(A_{\rm res}\setminus Z_\tau)\cup
  (B\cup\operatorname{redSupport}(\beta,J)).
\]
By the residual scan-stability input, which is the delete-and-reinsert
conclusion isolated in
\(\noderef{FixedSetHistoryCellRISIResidualScanStability}\) after the exact
read-set interface of
\(\noderef{FixedSetHistoryCellRISIRevealOrderScanConstruction}\), the scan
on \(A^\ast\) returns \(\tau\).

Now let \((\beta',\tau')\) be another realized pair whose forced-present
set is contained in \(A^\ast\), whose forced-absent set is disjoint from
\(A^\ast\), and whose complete blue trace agrees with \(A^\ast\) on every
blue coordinate of \(U\).  The full scan-recognition input applies to this
realized pair and the same assignment \(A^\ast\), so the same deterministic
scan also returns \(\tau'\).  A function has only one value at \(A^\ast\); hence
\(\tau'=\tau\).  This proves the first residual output.  Notice that branch
functionality and realized geometry alone would not give this conclusion:
the transcript equality is forced by the deterministic scan value.

For owner uniqueness, suppose the same residual assignment reconstructs
compatibly to two realized pairs \((\beta,\tau)\) and
\((\beta',\tau')\).  Apply the forward compatibility consequences to the
first reconstructed complete assignment.  Since compatibility contains
\(P_\tau\) in that reconstructed assignment and the fixed support
\[
B\cup\operatorname{redSupport}(\beta,J)
\]
is already inserted, every coordinate of
\[
P_\tau\setminus(B\cup\operatorname{redSupport}(\beta,J))
\]
must lie in the original residual assignment.  Since compatibility also
avoids \(A_\tau\setminus Z_\tau\) in the reconstructed assignment, and
coordinates outside \(Z_\tau\) are not deleted during reconstruction, the
original residual assignment is disjoint from
\(A_\tau\setminus Z_\tau\).  The same argument applied to
\((\beta',\tau')\) gives the residual span and avoidance conditions for the
second realized pair.

The hypotheses of the fixed residual first-mismatch/prefix-pruning input are
therefore satisfied: the same residual assignment lies in both residual
cylinders, and the two reconstructed assignments are compatible with their
respective realized cells.  This is exactly the uniqueness mechanism
formalized by
\(\noderef{FixedSetHistoryCellFixedBJResidualOverlapConsistency}\): a genuine
owner mismatch would produce a selected-present first mismatch, and the
prefix-pruning invariant rules out the corresponding compatibility
certificate.  Hence the two owners are equal, so
\(\beta=\beta'\) and \(\tau=\tau'\).
\end{proof}
